\documentclass[12pt]{ximera}
\typeout{Start loading xmPreamble.tex}%

% Add here extra macro's that are loaded automatically by all documents of claas 'ximera' or 'xourse' in this repo

%%
%%  Example:
%%
% \newcommand{\R}{\mathbb{R}}
\usepackage{tikz}
\usetikzlibrary{positioning, arrows.meta}
\usepackage{multicol}
\usepackage{enumitem}
\usepackage{caption}
\usepackage{amsmath}
\usepackage{amsthm}
\usepackage{fontenc}

%% ---------------------------------------------------------------
%% Shared worksheet macros.
%%
%% These came from the Summer 2026 worksheet set, where each worksheet carried its
%% own copy. They have to live here instead: a xourse inlines every activity into a
%% single document, so duplicate \newcommand definitions across activities collide,
%% and the xourse gobbles \input so a per-topic macro file never loads.
%%
%% They use \providecommand, NOT \newcommand. ximera.cls loads this file automatically,
%% and every activity also carries an explicit \typeout{Start loading xmPreamble.tex}%

% Add here extra macro's that are loaded automatically by all documents of claas 'ximera' or 'xourse' in this repo

%%
%%  Example:
%%
% \newcommand{\R}{\mathbb{R}}
\usepackage{tikz}
\usetikzlibrary{positioning, arrows.meta}
\usepackage{multicol}
\usepackage{enumitem}
\usepackage{caption}
\usepackage{amsmath}
\usepackage{amsthm}
\usepackage{fontenc}

%% ---------------------------------------------------------------
%% Shared worksheet macros.
%%
%% These came from the Summer 2026 worksheet set, where each worksheet carried its
%% own copy. They have to live here instead: a xourse inlines every activity into a
%% single document, so duplicate \newcommand definitions across activities collide,
%% and the xourse gobbles \input so a per-topic macro file never loads.
%%
%% They use \providecommand, NOT \newcommand. ximera.cls loads this file automatically,
%% and every activity also carries an explicit \typeout{Start loading xmPreamble.tex}%

% Add here extra macro's that are loaded automatically by all documents of claas 'ximera' or 'xourse' in this repo

%%
%%  Example:
%%
% \newcommand{\R}{\mathbb{R}}
\usepackage{tikz}
\usetikzlibrary{positioning, arrows.meta}
\usepackage{multicol}
\usepackage{enumitem}
\usepackage{caption}
\usepackage{amsmath}
\usepackage{amsthm}
\usepackage{fontenc}

%% ---------------------------------------------------------------
%% Shared worksheet macros.
%%
%% These came from the Summer 2026 worksheet set, where each worksheet carried its
%% own copy. They have to live here instead: a xourse inlines every activity into a
%% single document, so duplicate \newcommand definitions across activities collide,
%% and the xourse gobbles \input so a per-topic macro file never loads.
%%
%% They use \providecommand, NOT \newcommand. ximera.cls loads this file automatically,
%% and every activity also carries an explicit \input{../xmPreamble}, so this file is
%% read twice. That was harmless while it held only \usepackage lines (LaTeX skips a
%% package it has already loaded) but a second \newcommand is a hard error.
%%
%%   \blank[width]        fill-in-the-blank rule (default 2.5cm)
%%   \showwork            "show and explain your work" reminder
%%   \showworkline        ... plus which schedule line was used
%%   \showworkyear        ... plus which year's schedule was used
%%   \schedhead{...}      centred bold caption above a tiered-rate table
%%   \samesquares[scale]{left}{right}   two equal-area squares, labelled
%%   \samecubes[scale]{left}{right}     two equal-volume cubes, labelled
%% ---------------------------------------------------------------

\providecommand{\blank}[1][2.5cm]{\underline{\hspace{#1}}}
\providecommand{\showwork}{\textbf{REMEMBER TO SHOW AND EXPLAIN YOUR WORK.}}
\providecommand{\showworkline}{\textbf{REMEMBER TO SHOW AND EXPLAIN YOUR WORK.
  Explanation should include how and why you know which line of the
  schedule was used to calculate this amount.}}
\providecommand{\showworkyear}{\begin{sloppypar}\textbf{REMEMBER TO SHOW AND
  EXPLAIN YOUR WORK. Explanation should include how and why you know which
  line of the year's tax schedule was used to calculate this tax
  amount.}\end{sloppypar}}
\providecommand{\schedhead}[1]{\begin{center}\textbf{#1}\end{center}}

\providecommand{\samesquares}[3][1]{%
\begin{center}
{\footnotesize These squares are the \textbf{same size}.}

\vspace{1.5mm}
\begin{tikzpicture}[scale=#1,every node/.style={scale=#1}]
  \draw (0,0) rectangle (2.4,2.4);
  \node[anchor=north west] at (0.15,2.25) {Area:};
  \node[anchor=east]  at (-0.15,1.2) {#2};
  \node[anchor=north] at (1.2,-0.15) {#2};
  \begin{scope}[xshift=4.6cm]
    \draw (0,0) rectangle (2.4,2.4);
    \node[anchor=north west] at (0.15,2.25) {Area:};
    \node[anchor=east]  at (-0.15,1.2) {#3};
    \node[anchor=north] at (1.2,-0.15) {#3};
  \end{scope}
\end{tikzpicture}
\end{center}}

\providecommand{\samecubes}[3][1]{%
\begin{center}
{\footnotesize These cubes are the \textbf{same size}.}

\vspace{1.5mm}
\begin{tikzpicture}[scale=#1,every node/.style={scale=#1}]
  \foreach \dx in {0,5.2} {
    \begin{scope}[xshift=\dx cm]
      \draw (0,0) rectangle (2.0,2.0);
      \draw (0,2.0) -- (0.65,2.6) -- (2.65,2.6) -- (2.0,2.0);
      \fill[black!12] (2.0,0) -- (2.65,0.6) -- (2.65,2.6) -- (2.0,2.0) -- cycle;
      \draw (2.0,0) -- (2.65,0.6) -- (2.65,2.6) -- (2.0,2.0);
      \node[anchor=north west] at (0.1,1.9) {Volume:};
    \end{scope}
  }
  \node[anchor=east]  at (-0.15,1.0) {#2};
  \node[anchor=north] at (1.0,-0.15) {#2};
  \node[anchor=west]  at (2.25,0.4) {#2};
  \node[anchor=east]  at (5.05,1.0) {#3};
  \node[anchor=north] at (6.2,-0.15) {#3};
  \node[anchor=west]  at (7.45,0.4) {#3};
\end{tikzpicture}
\end{center}}

%% NOTE: do NOT add \usepackage to individual activity .tex files.
%% When a xourse inlines an activity it gobbles \documentclass and \input, but a bare
%% \usepackage still executes -- after \begin{document} -- and errors out.
%% All shared packages and macros belong here.
%%
%% ximera.cls already loads: amsmath, amssymb, amsthm, cancel, enumitem, graphicx,
%% hyperref, listings, pgfplots, tikz, url, xcolor. No need to re-load those.
, so this file is
%% read twice. That was harmless while it held only \usepackage lines (LaTeX skips a
%% package it has already loaded) but a second \newcommand is a hard error.
%%
%%   \blank[width]        fill-in-the-blank rule (default 2.5cm)
%%   \showwork            "show and explain your work" reminder
%%   \showworkline        ... plus which schedule line was used
%%   \showworkyear        ... plus which year's schedule was used
%%   \schedhead{...}      centred bold caption above a tiered-rate table
%%   \samesquares[scale]{left}{right}   two equal-area squares, labelled
%%   \samecubes[scale]{left}{right}     two equal-volume cubes, labelled
%% ---------------------------------------------------------------

\providecommand{\blank}[1][2.5cm]{\underline{\hspace{#1}}}
\providecommand{\showwork}{\textbf{REMEMBER TO SHOW AND EXPLAIN YOUR WORK.}}
\providecommand{\showworkline}{\textbf{REMEMBER TO SHOW AND EXPLAIN YOUR WORK.
  Explanation should include how and why you know which line of the
  schedule was used to calculate this amount.}}
\providecommand{\showworkyear}{\begin{sloppypar}\textbf{REMEMBER TO SHOW AND
  EXPLAIN YOUR WORK. Explanation should include how and why you know which
  line of the year's tax schedule was used to calculate this tax
  amount.}\end{sloppypar}}
\providecommand{\schedhead}[1]{\begin{center}\textbf{#1}\end{center}}

\providecommand{\samesquares}[3][1]{%
\begin{center}
{\footnotesize These squares are the \textbf{same size}.}

\vspace{1.5mm}
\begin{tikzpicture}[scale=#1,every node/.style={scale=#1}]
  \draw (0,0) rectangle (2.4,2.4);
  \node[anchor=north west] at (0.15,2.25) {Area:};
  \node[anchor=east]  at (-0.15,1.2) {#2};
  \node[anchor=north] at (1.2,-0.15) {#2};
  \begin{scope}[xshift=4.6cm]
    \draw (0,0) rectangle (2.4,2.4);
    \node[anchor=north west] at (0.15,2.25) {Area:};
    \node[anchor=east]  at (-0.15,1.2) {#3};
    \node[anchor=north] at (1.2,-0.15) {#3};
  \end{scope}
\end{tikzpicture}
\end{center}}

\providecommand{\samecubes}[3][1]{%
\begin{center}
{\footnotesize These cubes are the \textbf{same size}.}

\vspace{1.5mm}
\begin{tikzpicture}[scale=#1,every node/.style={scale=#1}]
  \foreach \dx in {0,5.2} {
    \begin{scope}[xshift=\dx cm]
      \draw (0,0) rectangle (2.0,2.0);
      \draw (0,2.0) -- (0.65,2.6) -- (2.65,2.6) -- (2.0,2.0);
      \fill[black!12] (2.0,0) -- (2.65,0.6) -- (2.65,2.6) -- (2.0,2.0) -- cycle;
      \draw (2.0,0) -- (2.65,0.6) -- (2.65,2.6) -- (2.0,2.0);
      \node[anchor=north west] at (0.1,1.9) {Volume:};
    \end{scope}
  }
  \node[anchor=east]  at (-0.15,1.0) {#2};
  \node[anchor=north] at (1.0,-0.15) {#2};
  \node[anchor=west]  at (2.25,0.4) {#2};
  \node[anchor=east]  at (5.05,1.0) {#3};
  \node[anchor=north] at (6.2,-0.15) {#3};
  \node[anchor=west]  at (7.45,0.4) {#3};
\end{tikzpicture}
\end{center}}

%% NOTE: do NOT add \usepackage to individual activity .tex files.
%% When a xourse inlines an activity it gobbles \documentclass and \input, but a bare
%% \usepackage still executes -- after \begin{document} -- and errors out.
%% All shared packages and macros belong here.
%%
%% ximera.cls already loads: amsmath, amssymb, amsthm, cancel, enumitem, graphicx,
%% hyperref, listings, pgfplots, tikz, url, xcolor. No need to re-load those.
, so this file is
%% read twice. That was harmless while it held only \usepackage lines (LaTeX skips a
%% package it has already loaded) but a second \newcommand is a hard error.
%%
%%   \blank[width]        fill-in-the-blank rule (default 2.5cm)
%%   \showwork            "show and explain your work" reminder
%%   \showworkline        ... plus which schedule line was used
%%   \showworkyear        ... plus which year's schedule was used
%%   \schedhead{...}      centred bold caption above a tiered-rate table
%%   \samesquares[scale]{left}{right}   two equal-area squares, labelled
%%   \samecubes[scale]{left}{right}     two equal-volume cubes, labelled
%% ---------------------------------------------------------------

\providecommand{\blank}[1][2.5cm]{\underline{\hspace{#1}}}
\providecommand{\showwork}{\textbf{REMEMBER TO SHOW AND EXPLAIN YOUR WORK.}}
\providecommand{\showworkline}{\textbf{REMEMBER TO SHOW AND EXPLAIN YOUR WORK.
  Explanation should include how and why you know which line of the
  schedule was used to calculate this amount.}}
\providecommand{\showworkyear}{\begin{sloppypar}\textbf{REMEMBER TO SHOW AND
  EXPLAIN YOUR WORK. Explanation should include how and why you know which
  line of the year's tax schedule was used to calculate this tax
  amount.}\end{sloppypar}}
\providecommand{\schedhead}[1]{\begin{center}\textbf{#1}\end{center}}

\providecommand{\samesquares}[3][1]{%
\begin{center}
{\footnotesize These squares are the \textbf{same size}.}

\vspace{1.5mm}
\begin{tikzpicture}[scale=#1,every node/.style={scale=#1}]
  \draw (0,0) rectangle (2.4,2.4);
  \node[anchor=north west] at (0.15,2.25) {Area:};
  \node[anchor=east]  at (-0.15,1.2) {#2};
  \node[anchor=north] at (1.2,-0.15) {#2};
  \begin{scope}[xshift=4.6cm]
    \draw (0,0) rectangle (2.4,2.4);
    \node[anchor=north west] at (0.15,2.25) {Area:};
    \node[anchor=east]  at (-0.15,1.2) {#3};
    \node[anchor=north] at (1.2,-0.15) {#3};
  \end{scope}
\end{tikzpicture}
\end{center}}

\providecommand{\samecubes}[3][1]{%
\begin{center}
{\footnotesize These cubes are the \textbf{same size}.}

\vspace{1.5mm}
\begin{tikzpicture}[scale=#1,every node/.style={scale=#1}]
  \foreach \dx in {0,5.2} {
    \begin{scope}[xshift=\dx cm]
      \draw (0,0) rectangle (2.0,2.0);
      \draw (0,2.0) -- (0.65,2.6) -- (2.65,2.6) -- (2.0,2.0);
      \fill[black!12] (2.0,0) -- (2.65,0.6) -- (2.65,2.6) -- (2.0,2.0) -- cycle;
      \draw (2.0,0) -- (2.65,0.6) -- (2.65,2.6) -- (2.0,2.0);
      \node[anchor=north west] at (0.1,1.9) {Volume:};
    \end{scope}
  }
  \node[anchor=east]  at (-0.15,1.0) {#2};
  \node[anchor=north] at (1.0,-0.15) {#2};
  \node[anchor=west]  at (2.25,0.4) {#2};
  \node[anchor=east]  at (5.05,1.0) {#3};
  \node[anchor=north] at (6.2,-0.15) {#3};
  \node[anchor=west]  at (7.45,0.4) {#3};
\end{tikzpicture}
\end{center}}

%% NOTE: do NOT add \usepackage to individual activity .tex files.
%% When a xourse inlines an activity it gobbles \documentclass and \input, but a bare
%% \usepackage still executes -- after \begin{document} -- and errors out.
%% All shared packages and macros belong here.
%%
%% ximera.cls already loads: amsmath, amssymb, amsthm, cancel, enumitem, graphicx,
%% hyperref, listings, pgfplots, tikz, url, xcolor. No need to re-load those.

\title{Other Tiered Systems}
\begin{document}
\begin{abstract}
Federal income tax is not the only thing structured in tiers. Comparing electricity
rates, Pell Grant contributions, and USPS shipping against the tax schedule reveals two
genuinely different designs --- one where a rate applies only to the amount above a
threshold, and one where crossing a threshold re-prices everything.
\end{abstract}
\maketitle

Federal income taxes are not the only things structured in this tiered way. Let's look
at some other contexts of tiered systems and see how they compare to the taxes system.

\section*{Electrical Use}

Below is a table of Columbus electricity rates.

\schedhead{DOP Electricity Rates, Effective May 1, 2025}

\begin{center}
\begin{tabular}{|l|p{5.6cm}|c|c|}
\hline
\textbf{Rate} & \textbf{Description} & \textbf{Customer charge} & \textbf{Energy per kWh} \\
\hline
\textbf{KW10} Residential Schedule A
& Applicable to residential users with \textbf{summer usage exceeding 700 kWh in any
month}.
& \$11.25 & \$0.11267 \\
\hline
\textbf{KW11} Residential Schedule A-1 (Small User)
& Applicable to residential users with \textbf{summer usage less than 700 kWh in any
month}.
& \$11.25 & \$0.09892 \\
\hline
\end{tabular}
\end{center}

{\footnotesize\url{https://www.columbus.gov/files/sharedassets/city/v/10/services/public-utilities/dop-rates.pdf}}

\begin{problem}
What is the maximum amount a ``small user'' can be charged in a month? Remember to
explain how and why you determined this.

About \$$\answer[tolerance=0.5]{80.40}$

\begin{freeResponse}
\end{freeResponse}

\begin{solution}
A small user is one whose usage stays \emph{below} 700 kWh, so the largest bill comes
from the largest qualifying usage --- just under 700 kWh. Taking 699 kWh:
\[
\$11.25 + \$0.09892 \cdot 699 = \$11.25 + \$69.15 = \$80.40.
\]
The \$11.25 customer charge is fixed and applies no matter how little you use, so the
bill is a flat fee plus a per-kWh charge.
\end{solution}
\end{problem}

\begin{problem}
What is the minimum amount a ``normal'' user can be charged in a month? Remember to
explain how and why you determined this.

About \$$\answer[tolerance=0.5]{90.23}$

\begin{freeResponse}
\end{freeResponse}

\begin{solution}
A normal user is one whose usage \emph{exceeds} 700 kWh, so the smallest bill comes from
the smallest qualifying usage --- just over 700 kWh. Taking 701 kWh:
\[
\$11.25 + \$0.11267 \cdot 701 = \$11.25 + \$78.98 = \$90.23.
\]
\end{solution}
\end{problem}

\begin{problem}
How do these minimum and maximum values compare? How is this different from the tax
brackets?

\begin{freeResponse}
\end{freeResponse}

\begin{solution}
The most a small user can be billed is about \$80.40, and the least a normal user can be
billed is about \$90.23. So there is a \textbf{gap of about \$9.84}: no residential
customer can receive a summer bill between those two amounts. The bill jumps
discontinuously as usage crosses 700 kWh.

This is fundamentally different from the tax brackets, and the difference is worth
stating precisely. In the tax schedule, crossing a bracket boundary changes the rate
applied to the \emph{additional} income only --- everything below the boundary keeps
being taxed the way it was. The result is continuous: an extra dollar of income never
causes a jump in tax owed.

Here, crossing 700 kWh re-prices \emph{every} kWh, from the first one, at the higher
rate of \$0.11267 instead of \$0.09892. Using one more kWh --- going from 699 to 701 ---
increases the bill by nearly \$10.

The practical consequence is real: a household at 699 kWh who turns on one more
appliance is worse off than before in a way that a taxpayer receiving a raise never is.
This kind of threshold is often called a \emph{cliff}, and it is exactly the thing people
mistakenly fear about tax brackets --- it just happens to be true here rather than there.
\end{solution}
\end{problem}

\begin{problem}
If someone was charged \$95.00 for electric use in a given month, are they a small user
or a normal user? How much electricity did they use?

About $\answer[tolerance=10]{743}$ kWh.

\begin{freeResponse}
\end{freeResponse}

\begin{solution}
Test both possibilities and see which one is self-consistent.

\emph{If they were a normal user}, then
\[
\$95.00 = \$11.25 + \$0.11267 \cdot k
\quad\Longrightarrow\quad
k = \frac{\$83.75}{\$0.11267} \approx 743.3 \text{ kWh}.
\]
That is above 700, which is exactly what being a normal user requires. Consistent.

\emph{If they were a small user}, then
\[
k = \frac{\$83.75}{\$0.09892} \approx 846.6 \text{ kWh},
\]
which is above 700 --- contradicting the requirement that a small user stay below 700.
So this case is impossible.

They are a normal user who used about 743 kWh. Notice the method: because the two rates
describe non-overlapping usage ranges, you assume one, solve, and check whether the
answer lands in the range that assumption requires.
\end{solution}
\end{problem}

\section*{Pell Grants}

Pell Grants are grants, not loans, meaning they do not have to be repaid. They are
awarded to undergraduate students to aid them in paying for college tuition. As with
taxes, we will make many assumptions to work past the various complicating factors. We
assume we are looking at a Pell Grant for a dependent student. After completing a
questionnaire about family income, it eventually comes to a table where an amount to be
paid by the parents is determined based on the Adjusted Available Income (AAI).

\begin{center}
\begin{tabular}{|l|l|}
\hline
\textbf{If the parents' AAI is} & \textbf{Then the parents' contribution from AAI is} \\
\hline
Less than $-$\$6,820      & $-$\$1,500 \\
\hline
$-$\$6,820 to \$20,600    & $22\%$ of AAI \\
\hline
\$20,601 to \$25,800      & \$4,532 plus $25\%$ of AAI over \$20,600 \\
\hline
\$25,801 to \$31,000      & \$5,832 plus $29\%$ of AAI over \$25,800 \\
\hline
\$31,001 to \$36,300      & \$7,340 plus $34\%$ of AAI over \$31,000 \\
\hline
\$36,301 to \$41,500      & \$9,142 plus $40\%$ of AAI over \$36,300 \\
\hline
\$41,501 or more          & \$11,222 plus $47\%$ of AAI over \$41,500 \\
\hline
\end{tabular}
\end{center}

\begin{problem}
If a parents' AAI is \$20,600, how much would their contribution be, according to this
table?

\$$\answer{4532}$

\begin{solution}
\$20,600 sits at the top of the second row, so the contribution is $22\%$ of AAI:
\[
0.22 \cdot \$20{,}600 = \$4{,}532.
\]
\end{solution}
\end{problem}

\begin{problem}
Is this the same set up as with taxes?

\begin{multipleChoice}
\choice[correct]{Yes --- each row charges a base amount plus a rate on only the amount above the threshold, exactly like the tax schedule.}
\choice{No --- crossing a threshold re-prices the entire AAI, like the electricity rates.}
\end{multipleChoice}

\begin{freeResponse}
\end{freeResponse}

\begin{solution}
Yes, this is the same structure as the tax schedule --- and you can verify it with the
same continuity check we used to find blanks A, B, and C.

The third row reads ``\$4,532 plus $25\%$ of AAI over \$20,600.'' That base of \$4,532 is
precisely the contribution we just computed at the top of the previous row. So the base
amount in each row is the total contribution accumulated from all the rows below it,
exactly as in the tax schedule.

Because of that, the contribution rises continuously: a family whose AAI increases by a
dollar never sees their expected contribution jump. Compare this with the electricity
rates, where crossing 700 kWh produced a \$9.84 discontinuity.

One small difference in presentation: this table uses whole-dollar boundaries (\$20,600
then \$20,601) rather than the tax schedule's ``over\ldots but not over'' phrasing.
That is a cosmetic difference in how the rows are written, not a difference in how the
system behaves.
\end{solution}
\end{problem}

\section*{Shipping Rates}

First Class shipping rates with USPS, by weight and zone:

\begin{center}
{\small\setlength{\tabcolsep}{3.5pt}
\begin{tabular}{|c|c|c|c|c|c|c|c|c|}
\hline
\textbf{Weight} & \textbf{Zones 1 \& 2} & \textbf{Zone 3} & \textbf{Zone 4} & \textbf{Zone 5}
& \textbf{Zone 6} & \textbf{Zone 7} & \textbf{Zone 8} & \textbf{Zone 9} \\
\hline
1--4 oz   & \$4.50 & \$4.60 & \$4.65 & \$4.70 & \$4.75 & \$4.80 & \$5.00 & \$5.00 \\
\hline
5--8 oz   & \$5.10 & \$5.15 & \$5.20 & \$5.25 & \$5.30 & \$5.40 & \$5.50 & \$5.50 \\
\hline
9--12 oz  & \$5.80 & \$5.85 & \$5.90 & \$5.95 & \$6.00 & \$6.15 & \$6.25 & \$6.25 \\
\hline
13 oz     & \$7.05 & \$7.10 & \$7.25 & \$7.30 & \$7.40 & \$7.55 & \$7.65 & \$7.65 \\
\hline
\end{tabular}}
\end{center}

{\footnotesize\url{https://www.usglobalmail.com/blog/usps-shipping-rates-by-weight-chart/}}

Zones are different distances from the origin address:

\begin{itemize}[nosep]
\item Zone 1: Local shipping (within 50 miles).
\item Zone 2: 51--150 miles.
\item Zone 3: 151--300 miles.
\item Zone 4: 301--600 miles.
\item Zone 5: 601--1000 miles.
\item Zone 6: 1001--1400 miles.
\item Zone 7: 1401--1800 miles.
\item Zone 8: 1801 miles or more.
\item Zone 9: US territories, freely associated states, and some military addresses.
\end{itemize}

\begin{problem}
Of the three systems in this activity --- electricity, Pell Grant contributions, and
USPS shipping --- which behaves most like the federal tax schedule, and which behaves
least like it? Explain what feature you are using to decide.

\begin{freeResponse}
\end{freeResponse}

\begin{solution}
The Pell Grant table behaves most like the tax schedule: a base amount plus a rate
applied only to the portion above a threshold, producing a continuous result.

USPS shipping behaves least like it. The rate does not depend on the amount above a
threshold at all --- every package from 1 to 4 oz costs the same \$4.50 to Zone 1,
and then the price steps to \$5.10 the moment it hits 5 oz. There is no ``rate per
ounce'' being applied to anything; the whole price is looked up from which block the
weight falls in. This is a \emph{step function}.

Electricity sits in between: like shipping, crossing the threshold re-prices everything
rather than just the excess, so it has a cliff. But unlike shipping, the charge within
each tier does still grow with usage at a per-kWh rate.

The feature that distinguishes them is what the threshold does to the amount
\emph{below} it. In the tax and Pell systems, income below a boundary keeps its old
treatment. In the electricity and shipping systems, crossing a boundary changes how the
entire quantity is priced --- which is what creates a jump.
\end{solution}
\end{problem}

\end{document}
